\section{Background}
\label{sec:background}

\subsection{LLM Inference Serving and Continuous Batching}
\label{subsec:serving}

Modern LLM serving engines process requests in \emph{iterations}: at each
decoding step, the engine advances every active sequence in the current batch
by one token, rather than running each request to completion before starting
the next. This scheme, known as continuous batching, allows new requests to
join a batch and finished requests to leave it between individual decoding
steps, so GPU compute is not left idle waiting for the longest sequence in a
static batch to finish. vLLM~\cite{kwon2023vllm}, the serving engine used in this work, combines
continuous batching with PagedAttention, a memory management scheme that
allocates the key-value cache in fixed-size, non-contiguous blocks analogous
to virtual memory pages. This avoids the internal and external fragmentation
that arises from pre-allocating a maximum-length contiguous cache per request,
which in turn allows more concurrent sequences to fit in a fixed amount of GPU
memory. Figure~\ref{fig:architecture_pipeline} illustrates where a scheduling
policy sits in this pipeline: it determines the order in which requests are
admitted into the continuous-batching engine, but it does not alter the engine
itself.

\begin{figure}[t]
    \centering
    \includegraphics[width=0.95\linewidth]{figures/architecture_pipeline.png}
    \caption{LLM serving pipeline with a pluggable scheduling policy. The
    policy re-orders incoming requests before they are admitted to vLLM's
    continuous-batching engine; the engine itself is unmodified.}
    \label{fig:architecture_pipeline}
\end{figure}

\subsection{Head-of-Line Blocking Under Unknown Output Length}
\label{subsec:hol}

Continuous batching removes idle GPU time between requests, but it does not
by itself protect short requests from being scheduled behind long ones. If the
admission order is FCFS, a request that happens to arrive early but will
generate hundreds of tokens occupies a batch slot for the full duration of its
decoding, and any request that arrives shortly after it but would have
finished in a few tokens must still wait its turn in the queue before it is
admitted. This is the LLM-serving analogue of HOL blocking in networking and
operating-systems scheduling: the request at the head of the queue determines
how long everything behind it waits, regardless of how short that following
work actually is. Because output length is not observable before decoding
completes, mitigating HOL blocking requires either estimating output length in
advance or using a proxy signal that correlates with it.

\subsection{Learning-to-Rank Scheduling}
\label{subsec:ltr-background}

One way to estimate output length in advance is to train a small model to
predict it, or a value that ranks requests similarly to it, directly from the
prompt text. The LTR approach this work builds on fine-tunes a compact
OPT-125M model~\cite{zhang2022opt} with a ranking loss (ListMLE) so that its output score orders
requests consistently with their true relative output lengths, rather than
predicting the exact token count. Requests are then admitted in order of this
learned score, giving requests expected to be short priority over requests
expected to be long. This is effective to the extent that the predictor's
training distribution matches production traffic, but the score is the
predictor's only signal: if the prompt distribution shifts or the predictor is
uncertain, there is no independent signal to fall back on. Prompt length is a
natural candidate for such a fallback signal, since it is known instantly and
without any additional inference cost; Figure~\ref{fig:composite_scoring}
shows how it is combined with the learned predictor score to form the
composite ranking used in this work, whose precise formulation is given in
Section~\ref{sec:methodology}. In practice, this ranking is computed offline over the current queue of waiting requests rather than per decoding step: each candidate prompt is tokenized and passed through the predictor once to obtain a scalar priority score, the queue is sorted by that score, and the resulting order is handed to the serving engine's continuous-batching scheduler, which then manages iteration-level admission as usual. The learned score therefore only ever influences which request is considered next, never how the engine executes it once admitted -- a separation that keeps the ranking step decoupled from, and safely composable with, whatever underlying batching or memory-management strategy the serving engine already uses. This also means the approach is not tied to a particular serving engine's internals: any system that exposes a pluggable admission order ahead of its own batching logic could, in principle, adopt the same ranking step without modification to its core execution path.
\begin{figure}[t]
    \centering
    \includegraphics[width=0.95\linewidth]{figures/composite_scoring_architecture.png}
    \caption{Composite scoring architecture. The learned LTR predictor score
    and a normalized prompt-length term are combined into a single ranking
    score used to order the admission queue.}
    \label{fig:composite_scoring}
\end{figure}